% !TEX root =./ECE343_Practice_main.tex
%
% L9 practice set (Impedance Matching II): admittance, the susceptance and
% distance at the matching point, and a parallel or stub match.

\fullwidth{\textbf{Learning Objective 3.4: } I can use the transmission line
equation to match a load to a source using a length of transmission line.
\\[4pt]\normalfont\small Show your work. Write a \textbf{Documentation} line at the
end (\textbf{None} if you did not collaborate). All lines are lossless, and $d$
is the distance from the load toward the source.}

\begin{questions}

%%%%%%%%%%%%%%%%%%%%%%%%%%%%%%%%%%%%%%%%%
\question \textbf{Admittance.}
\begin{parts}

	\part Substitute $z = 1/y$ into $\Gamma = \frac{z - 1}{z + 1}$ and show
	that $\Gamma = \frac{1 - y}{1 + y}$.
	\begin{solution}[1in]
	\eq{ \Gamma = \dfrac{1/y - 1}{1/y + 1} = \dfrac{1 - y}{1 + y} }
	after multiplying the numerator and denominator by $y$.
	\end{solution}

	\part Starting from the input impedance of a shorted stub and of an open
	stub of length $\ell_s$, find the normalized admittance of each.
	\begin{solution}[1.2in]
	\eq{
		\text{shorted:} \quad y &= \dfrac{\myz{0}}{j\myz{0}\tan\beta\ell_s} = -j\cot\beta\ell_s \\
		\text{open:} \quad y &= \dfrac{\myz{0}}{-j\myz{0}\cot\beta\ell_s} = j\tan\beta\ell_s
	}
	\end{solution}

	\part Explain in one sentence why a matching component that connects
	across the line is easier to design with admittance than with impedance.
	\begin{solution}[0.7in]
	Admittances in parallel add, so the component's admittance simply adds to
	the line's admittance at that point.
	\end{solution}

\end{parts}

\newpage
%%%%%%%%%%%%%%%%%%%%%%%%%%%%%%%%%%%%%%%%%
\question \textbf{A parallel match.} The $100\ \ohms$ load from the Lesson 8
practice set terminates a $50\ \ohms$ line at $2\ \ghz$, with
$\lambda = 10\ \text{cm}$ in the line. Its reflection coefficient is
$\Gamma = 0.333\angle 0^\circ$.
\begin{parts}

	\begin{minipage}[t]{0.58\linewidth}
		\part Find $y_L$ and the two values of $b$ at the matching points.
		\begin{solution}[1in]
		\eq{
			y_L &= \dfrac{50}{100} = 0.5 \\
			b &= \pm\dfrac{2 \times 0.333}{\sqrt{1 - 0.111}} = \pm 0.707
		}
		\end{solution}
	\end{minipage}\hfill%
	\begin{minipage}[t]{0.37\linewidth}
		\raggedleft \vspace{12pt}
		$y_L$ = \ansbox[1.3in]{$0.5$}

		\vspace{4pt}
		$b$ = \ansbox[1.3in]{$\pm 0.707$}
	\end{minipage}

	\begin{minipage}[t]{0.58\linewidth}
		\part Find the phase $\theta_d$ and the distance $d$ for each
		matching point.
		\begin{solution}[1.8in]
		$\tan^{-1}(0.707/2) = 19.5^\circ$, so
		$\theta_d = \mp\lp 90^\circ + 19.5^\circ\rp = \mp 109.5^\circ$.
		\eq{
			b = +0.707: \quad d &= \dfrac{0^\circ + 109.5^\circ}{720^\circ}\lambda \\
			&= 0.152\lambda = 1.52\ \text{cm} \\
			b = -0.707: \quad d &= \dfrac{0^\circ - 109.5^\circ + 360^\circ}{720^\circ}\lambda \\
			&= 0.348\lambda = 3.48\ \text{cm}
		}
		\end{solution}
	\end{minipage}\hfill%
	\begin{minipage}[t]{0.37\linewidth}
		\raggedleft \vspace{12pt}
		$d_1$ = \ansbox[1.3in]{$0.152\lambda$}

		\vspace{4pt}
		$d_2$ = \ansbox[1.3in]{$0.348\lambda$}
	\end{minipage}

	\begin{minipage}[t]{0.58\linewidth}
		\part Find the parallel component that completes each match, and its
		value.
		\begin{solution}[1.7in]
		The susceptance to cancel is $0.707/50 = 14.1\ \text{mS}$.
		At $d_1$, $y = 1 + j0.707$, so we add a parallel inductor:
		\eq{ L = \dfrac{1}{2\pi \times 2\times10^{9} \times 0.0141} = 5.63\ \text{nH} }
		At $d_2$, $y = 1 - j0.707$, so we add a parallel capacitor:
		\eq{ C = \dfrac{0.0141}{2\pi \times 2\times10^{9}} = 1.13\ \text{pF} }
		\end{solution}
	\end{minipage}\hfill%
	\begin{minipage}[t]{0.37\linewidth}
		\raggedleft \vspace{12pt}
		at $d_1$: \ansbox[1.3in]{$5.63\ \text{nH}$}

		\vspace{4pt}
		at $d_2$: \ansbox[1.3in]{$1.13\ \text{pF}$}
	\end{minipage}

\end{parts}

\newpage
%%%%%%%%%%%%%%%%%%%%%%%%%%%%%%%%%%%%%%%%%
\question \textbf{Stubs in place of components.} Replace each component from
Question 2 with a stub of the same $50\ \ohms$ line.
\begin{parts}

	\begin{minipage}[t]{0.58\linewidth}
		\part At $d_1$ the stub must supply $b_s = -0.707$. Find the shortest
		shorted stub and the shortest open stub that do.
		\begin{solution}[1.5in]
		\eq{
			\text{shorted:} \quad \cot\beta\ell_s &= 0.707 \\
			\beta\ell_s &= 54.7^\circ, \quad \ell_s = 0.152\lambda \\
			\text{open:} \quad \tan\beta\ell_s &= -0.707 \\
			\beta\ell_s &= 144.7^\circ, \quad \ell_s = 0.402\lambda
		}
		\end{solution}
	\end{minipage}\hfill%
	\begin{minipage}[t]{0.37\linewidth}
		\raggedleft \vspace{12pt}
		shorted = \ansbox[1.3in]{$0.152\lambda$}

		\vspace{4pt}
		open = \ansbox[1.3in]{$0.402\lambda$}
	\end{minipage}

	\begin{minipage}[t]{0.58\linewidth}
		\part At $d_2$ the stub must supply $b_s = +0.707$. Find the shortest
		shorted stub and the shortest open stub that do.
		\begin{solution}[1.5in]
		\eq{
			\text{shorted:} \quad \cot\beta\ell_s &= -0.707 \\
			\beta\ell_s &= 125.3^\circ, \quad \ell_s = 0.348\lambda \\
			\text{open:} \quad \tan\beta\ell_s &= 0.707 \\
			\beta\ell_s &= 35.3^\circ, \quad \ell_s = 0.098\lambda
		}
		\end{solution}
	\end{minipage}\hfill%
	\begin{minipage}[t]{0.37\linewidth}
		\raggedleft \vspace{12pt}
		shorted = \ansbox[1.3in]{$0.348\lambda$}

		\vspace{4pt}
		open = \ansbox[1.3in]{$0.098\lambda$}
	\end{minipage}

	\part Which line-and-stub design uses the least total line? Give its two
	lengths in centimeters.
	\begin{solution}[0.9in]
	The line of $0.152\lambda$ with a shorted stub of $0.152\lambda$, a total of
	$0.304\lambda$: $1.52\ \text{cm}$ of line and a $1.52\ \text{cm}$ stub. The
	other short choice, $0.348\lambda$ with an open stub of $0.098\lambda$, totals
	$0.446\lambda$.
	\end{solution}

\end{parts}

\newpage
%%%%%%%%%%%%%%%%%%%%%%%%%%%%%%%%%%%%%%%%%
\question \textbf{A match at the load.} A load $\myz{L} = 25 - j25\ \ohms$
terminates a $50\ \ohms$ line at $2\ \ghz$.
\begin{parts}

	\begin{minipage}[t]{0.58\linewidth}
		\part Find $y_L$.
		\begin{solution}[1in]
		\eq{ y_L = \dfrac{1}{0.5 - j0.5} = \dfrac{0.5 + j0.5}{0.5} = 1 + j1 }
		\end{solution}
	\end{minipage}\hfill%
	\begin{minipage}[t]{0.37\linewidth}
		\raggedleft \vspace{12pt}
		$y_L$ = \ansbox[1.3in]{$1 + j1$}
	\end{minipage}

	\begin{minipage}[t]{0.58\linewidth}
		\part The conductance is already 1, so $d = 0$ is a matching point.
		Find the single parallel component, connected across the load, that
		matches it.
		\begin{solution}[1.2in]
		It must supply $-j1$, or $-j20\ \text{mS}$, so it is an inductor:
		\eq{ L = \dfrac{1}{2\pi \times 2\times10^{9} \times 0.02} = 3.98\ \text{nH} }
		\end{solution}
	\end{minipage}\hfill%
	\begin{minipage}[t]{0.37\linewidth}
		\raggedleft \vspace{12pt}
		$L$ = \ansbox[1.3in]{$3.98\ \text{nH}$}
	\end{minipage}

\end{parts}

\vspace{1.5em}
\noindent\textbf{Documentation:} \rule[-2pt]{0.6\linewidth}{0.4pt}

\end{questions}
